\subsection{The administration service}
\label{sec:ha:disabled:admin}

The administration service is the only component permitted to change credential
state. Members never contact it. Authentication services receive what it sends
them.

\textbf{One instance runs in both configurations, by decision.} The service is
not on the path an order takes, so its absence stops no member trading or logging
on. Credentials change rarely, so the venue tolerates its absence for as long as
a restart takes. A second instance would need arbitration of its own to remain
the single writer, which is a large mechanism for a component that is idle for
most of a trading day.

That decision carries one condition: everything the single instance must retain
has to survive the instance.

\subsubsection{What state it holds}
\label{sec:ha:disabled:admin:state}

\textbf{The credential state.} The service writes this state, and the database
holds the original.

\textbf{Which authentication services have applied which change.} No other
component holds this. An authentication service holds no record of a change it
never received, and the database records what the credentials are rather than
which instance has been told.

\subsubsection{Where the copy that outlives it lives}
\label{sec:ha:disabled:admin:durable}

The database holds the credential state. Nothing holds the record of which
instances have applied which change.

\begin{gap}{}
  Which instances have applied a change is held only in the running process. An
  administration service that restarts between making a change and distributing
  it to every instance cannot afterwards establish which instances were missed.
\end{gap}

\subsubsection{How a restarted instance becomes current}
\label{sec:ha:disabled:admin:recovery}

The service reads the credential state from the database, and that state is
complete. It cannot recover the record of what it had distributed, because
nothing wrote that record down.

The service has one such operation rather than two, as the gateway in
\cref{sec:ha:disabled:gateway} and the authentication service in
\cref{sec:ha:disabled:auth} do. An administration service is never promoted,
because the venue never runs a second one to promote.

\subsubsection{What the member is told}
\label{sec:ha:disabled:admin:member}

Nothing directly. An authentication service answers a member's logon, and the
requirements in this subsection exist so that the answer does not depend on an
administration service running at the time.

\begin{requirement}{R-0096}{A credential change is recorded before it is distributed}
  A change shall be written to the source of truth before it is sent to any
  instance, so that a failure between the two loses the distribution rather than
  the change.
  \rationale{The next instance to reload from the database undoes a change that was
  distributed but not recorded, and nothing reports that. A change that was
  recorded but not distributed appears as a difference between instances, and the
  venue can correct it.}
  \uncovered{no scenario interrupts a credential change}
\end{requirement}

\begin{requirement}{R-0097}{The record of which instances are behind outlives the service that made it}
  Where an instance did not receive a change, that shall be recorded durably, and
  shall be acted on by whichever administration service is running when the
  instance next becomes reachable.
  \rationale{\reqref{R-0046} accepts a change while an instance is unavailable, on the basis
  that the venue applies it later. Where only a running process holds the
  obligation to apply that change, the obligation is lost whenever the
  administration service restarts before the instance returns. The credential state
  is then permanently inconsistent, and nothing records that it is.}
  \uncovered{no scenario restarts the administration service}
\end{requirement}

\begin{requirement}{R-0112}{The venue trades while the store it depends on is unreachable}
  No component shall require a connection to the durable store in order to admit
  a member, accept an order, or answer one. What it needs from the store shall be
  read when it starts and supplied to it thereafter.
  \rationale{Where the venue needs the store to be reachable in order to work, every outage of
  the store becomes a trading outage, and the venue takes on a failure whose
  likelihood has nothing to do with anything the venue does. Admission already
  works this way: an authentication service reads an export when it starts and is
  sent each change afterwards. Orders that outlive the day are to work the same
  way.}
  \uncovered{no scenario runs the venue with the store unreachable}
\end{requirement}

\begin{requirement}{R-0098}{Members trade and log on while no administration service is running}
  The absence of an administration service shall not prevent a member logging on,
  placing orders, or being answered.
  \rationale{The administration service writes credential state, and answers nothing when a
  member is admitted. Making a member's logon depend on it would place a component
  that is idle for most of a trading day on the path of every session that starts.}
  \uncovered{no scenario runs the venue without an administration service}
\end{requirement}
